 \chapter*{\textbf{Abstract}}
The agricultural sector faces a critical labor shortage, particularly in the harvesting of delicate crops like green asparagus, which requires selective picking and labor-intensive manual handling. This thesis presents the design, development, and validation of an autonomous harvesting framework aimed at addressing this challenge. As a proof of concept, the system integrates a custom-built planar robotic arm which serves as a testbed for evaluating the robustness and precision of the computer vision and control algorithms.

The system's core intelligence relies on a robust computer vision pipeline that 
combines deep learning with 3D spatial analysis. A Mask R-CNN model is employed 
for instance segmentation to identify asparagus spears from RGB imagery, while 
3D data is projected onto 2D space to establish a unified spatial reference frame that streamlines computational workflows. Leveraging the segmentation 
masks generated from the RGB imagery, the system achieves precise isolation and localization of the asparagus 
targets. These 2D detections are subsequently transformed into 3D spatial coordinates, facilitating accurate 
kinematic targeting for the robotic manipulator to execute its reach and operation tasks.

A hybrid control architecture coordinates the operation, with high-level Python 
software performing inverse kinematics and path planning, and distributed 
Arduino microcontrollers executing real-time actuator control via low-level 
firmware. Experimental validation, conducted in a controlled laboratory 
environment with the planar manipulator, confirmed the effectiveness of the 
vision system. The results demonstrate the system's ability to consistently 
identify and reach target locations, validating the proposed computational 
framework as a foundation for future field-deployable agricultural robots.

